\documentclass[10pt,a4paper]{amsart}
\usepackage[margin=19mm]{geometry}
\usepackage{amsmath,amssymb,mathtools}
\usepackage[scr=rsfs,cal=euler]{mathalfa}
\usepackage{xfrac}
\usepackage[unicode,hidelinks,bookmarks=false]{hyperref}
\newcommand{\ie}{i.e.\ }
\newcommand{\cf}{cf.\ }
\newcommand{\resp}{respectively\ }
\newcommand{\Subsection}{\subsection}
\providecommand{\nobreakdash}{\nobreak\hbox{-}}
\setlength{\emergencystretch}{2em}
\multlinegap0pt
\usepackage{bm}
\usepackage{bbm}
\usepackage{dsfont}
\usepackage{xspace}
\usepackage{cancel}
\usepackage{mathtools}
\usepackage{amsthm}
\makeatletter
\@ifundefined{theorem}{\newtheorem{theorem}{Theorem}}{}
\makeatother
\title[The Jones polynomial in systems with Periodic Boundary Condi]{The Jones polynomial in systems with Periodic Boundary Conditions}
\author{Kasturi Barkataki, Eleni Panagiotou}
\date{Source: arXiv:2309.14572v1 (CC BY 4.0)}
\begin{document}
\maketitle

\noindent\emph{Source and licence.} The Jones polynomial in systems with Periodic Boundary Conditions. Version arXiv:2309.14572v1 (CC BY 4.0), distributed under the Creative Commons Attribution 4.0 International licence, as recorded in the arXiv licence metadata for this exact version and shown on the abstract page https://arxiv.org/abs/2309.14572v1. Full rights notice: licenses/arxiv-2309.14572-CC-BY-4.0.txt.\par\smallskip

\noindent\textit{Editorial scope.} Counted unit: the Theorem whose source label is \texttt{VHN-exp} in the article (source0.tex lines 498--509, source label \texttt{VHN-exp}), delivered together with the article's own complete original proof of it (source0.tex lines 510--541, one \texttt{proof} environment of 32 lines). No mathematics is written, repaired or re-derived: statement and proof are the article's own text line for line. Required context retained uncounted: none. Every definition, display and label of those lines is retained unchanged. Omitted from this package and not redistributed: 1--497, 542--567. Nothing inside a retained excerpt is omitted or altered: every retained body is delivered exactly as the article's lines are written, so no omission receipt of any kind accompanies this package. Citations: the retained excerpts carry no citation command at all, so no bibliography entry of the article is transcribed and no reference is redistributed. Reference adaptation: none was needed and none was made; every reference inside the delivered unit resolves inside it, so no unresolved reference or citation is produced. Printed slips kept literal: no printed word, symbol, hypothesis or number of the counted statement or of its proof was altered. Layout and class: the article's own class is \texttt{article} with packages including amsmath, amssymb, amsthm, xcolor, algorithm2e, setspace, xr-hyper, hyperref, fullpage, graphicx, calc; the wrapper is the lane's pinned \texttt{amsart} editorial wrapper at 10pt on A4, which restates the theorem-like environments the retained text needs and adds the article's own macro definitions that the retained text needs (none), with extra wrapper packages (bm, bbm, dsfont, xspace, cancel, mathtools). The delivered numbering is the unit's own. Assets: no retained span contains a figure environment, a graphics-inclusion command, a PDF-page inclusion, a table or a TikZ picture; no image file is copied into this package, the package declares no asset dependency and no image file is redistributed with it. Licence: version arXiv:2309.14572v1 of this article is distributed under the Creative Commons Attribution 4.0 International licence, as recorded in the arXiv licence metadata for this exact version and shown on the abstract page https://arxiv.org/abs/2309.14572v1. A full rights notice is delivered at licenses/arxiv-2309.14572-CC-BY-4.0.txt. Characters: every retained excerpt is pure ASCII, so the delivered engine, XeLaTeX with its default Latin Modern text font, reports no missing character.
\begin{theorem} (Jones polynomial for any cutoff of a system with 1 closed chain in 1 PBC)
Let $C$ be the base cell of a periodic system generated by 1 closed chain in 1 PBC and let $\mathcal{MU}_C$ be its minimal collective unfolding consisting of $|\mathcal{MU}_C|$ copies of $C$ and $\mathcal{L}_C$ its minimal periodic link. Let $\mathcal{N}_{\mathcal{MU}_C}$ be $(2N-1)( |\mathcal{MU}_C|) - (N-1)$ repeated copies of the base cell $C$ (along the PBC faces) and let $(\mathcal{L}_{C})_N$, denote the link composed by all images of $I$ intersecting or contained in $\mathcal{N}_{\mathcal{MU}_C}$.  The Jones polynomial of $(\mathcal{L}_C)_{N}$ can be expressed as
\begin{equation}
    \begin{split}\displaystyle
        \mathsf{V}\left((\mathcal{L}_C)_{N}\right)  &=(-A)^{-(N-1)\mathsf{SLK}_P(\mathcal{L}_C)} d^{N-1} {\mathsf{V}_P(C)}^N + \tilde{\Lambda}\\
    \end{split}
    \label{VHN-exp}
\end{equation}
\noindent where ${\mathsf{V}_P(C)}$ is the Periodic Jones polynomial of $C$, $\mathsf{SLK}_P(\mathcal{L}_C)$ denotes the periodic linking with self images of copies of $\mathcal{L}_C$. The term $(-A)^{-(N-1)\mathsf{SLK}_P(\mathcal{L}_C)} d^{N-1} {\mathsf{V}_P(C)}^N$ is the contribution of one of the states of the Jones polynomial of $(\mathcal{L}_C)_{N}$ which leads to disconnection of all copies of $\mathcal{L}_{C}$ in $(\mathcal{L}_{C})_N$ and the remainder term, $\tilde{\Lambda}$, is contributed by the states which do not lead to disconnection of the $N$ 
 copies of $\mathcal{L}_{C}$ in $(\mathcal{L}_{C})_N$.
\label{corr1PBCjones}
\end{theorem}

\begin{proof}
Notice that in $\mathcal{N}_{\mathcal{MU}_C}$, which consists of $(2N-1)( |\mathcal{MU}_C|) - (N-1)$ cells, we can identify exactly $N$ minimal unfoldings (consisting of $|\mathcal{MU}_C|$ cells each) separated by $|\mathcal{MU}_C|-1$ cells (Thus we have $(N-1)(|\mathcal{MU}_C|-1)$ intermediate cells). Each minimal unfolding defines a minimal periodic link, $\displaystyle \mathcal{L}_C$, and, due to the separation between these minimal unfoldings, these minimal periodic links have no common components. Indeed, a component that intersects $\mathcal{MU}_C$ can intersect at most $|\mathcal{MU}_C|-1$ cells not in $\mathcal{MU}_C$. Moreover, the components that belong in these minimal periodic links are exactly all the components intersecting or inside $\mathcal{N}_{\mathcal{MU}_C}$, since any component in the $|\mathcal{MU}_C|-1$ intermediate cells must intersect one of the neighboring $\mathcal{MU}_C$.

Let $\vec{\xi}\in S^2$ and project the periodic system to the plane with normal vector $\vec{\xi}$. Let $c$ denote the crossings in ${(\mathcal{L}_C)}_{\vec{\xi}}$, we call them self-crossings, and let $k$ be the crossings between any consecutive pair of disjoint (in terms of components) copies of $\displaystyle {(\mathcal{L}_C)}_{\vec{\xi}}$, we call them shared crossings.  Then
 $\displaystyle \left \langle {((\mathcal{L}_C)_{N})}_{\vec{\xi}} \right \rangle$ can be expanded as, $\displaystyle \sum_{\mathcal{T}} \langle 
\mathcal{T} \rangle$, where $\mathcal{T}$ is any state obtained by resolving the $(N-1)k$ shared crossings in $\displaystyle {((\mathcal{L}_{C})_{N})}_{\vec{\xi}}$. Any such $\mathcal{T}$ is in fact a knot/link with a total of $\displaystyle Nc$ unresolved crossings.  Among the $2^{(N-1)k}$ such states, there exists a unique state, $\mathcal{S}$, which leads to the disconnection of  the $N$ copies of ${(\mathcal{L}_C)}_{\vec{\xi}}$ in  $\displaystyle {((\mathcal{L}_C)_{N})}_{\vec{\xi}}$. Smoothing along the shared crossings between any two copies of ${(\mathcal{L}_C)}_{\vec{\xi}}$ by respecting orientation, contributes a factor $A$, if the crossing is positive or a factor $A^{-1}$, if the crossing is negative. Thus, disconnecting a copy of  ${(\mathcal{L}_C)}_{\vec{\xi}}$ accumulates a factor $A^\theta$, where $\theta= 2\mathsf{SLK}_P({(\mathcal{L}_C)}_{\vec{\xi}})= 2\sum_{\vec{v}} \mathsf{L}({(\mathcal{L}_C)}_{\vec{\xi}}, {(\mathcal{L}_C)}_{\vec{\xi}} + \vec{v})$ (where $\mathsf{L}$ denotes the half algerbaic sum of shared crossings) the . Therefore, $A^{(N-1)2\mathsf{SLK}_P({(\mathcal{L}_C)}_{\vec{\xi}})}$ is the net smoothing factor associated with the state, $\mathcal{S}$. Then,
\begin{equation}
    \begin{split}\displaystyle
        \left\langle {((\mathcal{L}_C)_{N})}_{\vec{\xi}} \right\rangle  &= A^{2(N-1)\mathsf{SLK}_P({(\mathcal{L}_C)}_{\vec{\xi}})} d^{N-1} \left\langle {(\mathcal{L}_C)}_{\vec{\xi}} \right\rangle^N + \Lambda,
    \end{split}
    \label{eq-theta-lam}
\end{equation}

\noindent where $\Lambda$ is the contribution to the Jones polynomial from the remaining $2^{(N-1)k}-1$ states.

The Jones polynomial is obtained upon the normalisation of Equation \ref{eq-theta-lam} by the diagrammatic writhe of $ {((\mathcal{L}_C)_{(N)})}_{\vec{\xi}}$ :
\begin{equation}
\begin{split}\displaystyle \mathsf{Wr}\left({((\mathcal{L}_C)_{(N)})}_{\vec{\xi}}\right) &= \sum_{i=1}^{N} \mathsf{Wr}\left((\mathcal{L}_C^i)_{\vec{\xi}}\right) + \sum_{i,j=1, i\neq j}^{N-1} \mathsf{L}\left((\mathcal{L}_C^i)_{\vec{\xi}},(\mathcal{L}_C^{j})_{\vec{\xi}}\right)\\&= N \mathsf{Wr}\left({(\mathcal{L}_C)}_{\vec{\xi}}\right) + (N-1) \mathsf{SLK}_P\left({(\mathcal{L}_C)}_{\vec{\xi}}\right),
\end{split}
\label{wrr1}
\end{equation}
where $\mathsf{Wr}\left({(\mathcal{L}_C^i)}_{\vec{\xi}}\right)$ is the diagrammatic writhe of a copy of ${(\mathcal{L}_C)}_{\vec{\xi}}$; $\mathsf{L}\left((\mathcal{L}_C^i)_{\vec{\xi}},(\mathcal{L}_C^{j})_{\vec{\xi}}\right)$ is the diagrammatic linking number between the $i^{th}$ and $j^{th}$ copies of ${(\mathcal{L}_C)}_{\vec{\xi}}$ and $\mathsf{SLK}_P\left({(\mathcal{L}_C)}_{\vec{\xi}}\right)$ is the periodic linking  with self images of ${(\mathcal{L}_C)}_{\vec{\xi}}$ (accounting only shared crossings). The expression for the Jones polynomial of $(\mathcal{L}_C)_{N}$ is then 

\begin{equation}
    \begin{split}\displaystyle
        \mathsf{V}\left((\mathcal{L}_C)_{N}\right) &= {(-A^3)}^{-\mathsf{Wr}\left({((\mathcal{L}_C)_{N})}_{\vec{\xi}}\right)} \left\langle {((\mathcal{L}_C)_{N})}_{\vec{\xi}} \right\rangle\\
         &=(-A)^{-(N-1)\mathsf{SLK}_P\left({(\mathcal{L}_C)}_{\vec{\xi}}\right)} d^{N-1} {\mathsf{V}\left({(\mathcal{L}_C)}_{\vec{\xi}}\right)}^N + \tilde{\Lambda}\\
        &=(-A)^{-(N-1)\mathsf{SLK}_P\left({(\mathcal{L}_C)}\right)} d^{N-1} {\mathsf{V}\left({(\mathcal{L}_C)}\right)}^N + \tilde{\Lambda},\\
    \end{split}
\end{equation}
where $\displaystyle \tilde{\Lambda} = \Lambda {(-A^3)}^{-\mathsf{Wr}\left({((\mathcal{L}_C)_{N})}_{\vec{\xi}}\right)}$. 
\end{proof}

\end{document}
